% sections/01-analisis-trafico.tex — punto 5.1 de la guía
\section{Análisis de tráfico con Wireshark}

\subsection{Puesta en marcha de la captura}

Wireshark se encuentra disponible en Kali Linux dentro del menú de aplicaciones (categoría \emph{09 -- Sniffing \& Spoofing}) y también puede lanzarse desde la terminal, teniendo en cuenta que la captura requiere privilegios sobre la interfaz de red:

\begin{lstlisting}[language=bash]
sudo wireshark &
\end{lstlisting}

Una vez abierta la aplicación se ingresó al botón \emph{Capture Options}, desde el cual se seleccionó la interfaz sobre la que se realizaría la escucha (la interfaz cableada \texttt{eth0} o la inalámbrica \texttt{wlan0}, según el montaje), dejando activo el modo promiscuo para que la tarjeta entregue al sistema todas las tramas que observa y no solo las dirigidas a su propia dirección MAC:

\captura{QWE-535.png}{Ventana \emph{Capture Options} de Wireshark con la interfaz de red seleccionada para la escucha.}

\begin{analisis}




En la ventana Capture Options se listan las interfaces disponibles; se seleccionó eth0, por donde circula el tráfico de interés, y se mantuvo activa la casilla Enable promiscuous mode on all interfaces. El modo promiscuo hace que la tarjeta entregue al sistema todas las tramas que llegan al medio y no solo las dirigidas a su dirección MAC; sin embargo, es una condición necesaria pero no suficiente, dado que en una red conmutada el switch solo reenvía al puerto el tráfico propio y el de difusión, por lo que capturar conversaciones de terceros exige además un puerto espejo (SPAN) o una posición intermedia en la ruta. En redes inalámbricas el equivalente es el modo monitor, que aparece como opción independiente en la misma ventana. La imagen se tomó al reabrir la ventana con la captura ya en curso, para documentar la interfaz utilizada.
\end{analisis}

Posterior a esto Wireshark comenzó a capturar los paquetes que pasan por la interfaz seleccionada, mostrándolos en tiempo real en el panel de lista de paquetes:

\captura{wireshark-captura-viva.png}{Captura en curso sobre la interfaz seleccionada.}

\begin{analisis}


La lista de paquetes se puebla en tiempo real con el tráfico de la interfaz. En la figura, filtrada con ip.addr == 10.38.90.239 para conservar solo el tráfico del equipo de pruebas, conviven consultas DNS (como la resolución de www.gstatic.com), el establecimiento de conexiones TCP y sesiones TLS 1.3 hacia servicios externos. En el paquete 404 se distingue un Client Hello cuya extensión SNI revela el nombre del sitio, seguido de registros Application Data cifrados, lo que anticipa la diferencia de visibilidad entre el tráfico en claro y el cifrado que se analiza a continuación.
\end{analisis}

\subsection{Credenciales sobre HTTP}

Se realizó la búsqueda de una página de inicio de sesión que hiciera uso de HTTP y se ingresaron unas credenciales de prueba en los campos de usuario y contraseña (nunca credenciales reales, dado que el tráfico queda registrado en el archivo de captura). Para aislar el envío del formulario dentro de todo el tráfico capturado se aplicaron los siguientes filtros de visualización:

\begin{lstlisting}
http.request.method == "POST"
urlencoded-form
\end{lstlisting}

\captura{http-post-credenciales.png}{Paquete HTTP POST con los campos del formulario de autenticación en texto claro.}

\begin{analisis}


Tras aplicar el filtro queda aislada la petición POST /submit/ enviada al sitio de pruebas http-login.badssl.com. Al desplegar en su árbol de disección el nodo HTML Form URL Encoded se leen directamente los pares nombre-valor del formulario (usuario y contraseña tal como se escribieron), sin ninguna operación de descifrado ni herramienta adicional, lo que evidencia que HTTP no ofrece garantía alguna de confidencialidad.
\end{analisis}

\pregunta{Al analizar los paquetes HTTP capturados, ¿se evidencia la información de las credenciales? Justifique su respuesta.}
\begin{respuesta}
Sí se evidencia la información de las credenciales, y la justificación está en el propio diseño del protocolo: HTTP transporta la petición y su cuerpo en texto claro sobre TCP, sin ninguna capa de cifrado intermedia, de forma que cuando el formulario de autenticación se envía mediante el método POST los campos de usuario y contraseña viajan como pares nombre-valor codificados en \eng{url-encoding} dentro del cuerpo de la petición (\texttt{application/x-www-form-urlencoded}), y cualquier equipo ubicado en la ruta del paquete, o con acceso al mismo segmento mediante un puerto espejo o un ataque de intermediario, puede leerlos sin realizar operación alguna de descifrado. En Wireshark esto se observa directamente en el árbol de disección del paquete, bajo el nodo \emph{HTML Form URL Encoded}, o de forma más cómoda siguiendo el flujo completo con \emph{Follow $\rightarrow$ HTTP Stream}, donde se reconstruye la conversación entre el navegador y el servidor tal como se transmitió. Cabe mencionar que la exposición no se limita a las credenciales: también quedan legibles las \eng{cookies} de sesión, los encabezados y el contenido devuelto por el servidor, de manera que un atacante podría suplantar la sesión incluso sin conocer la contraseña.
\end{respuesta}

\subsection{Credenciales sobre HTTPS}

Posteriormente se repitió el mismo ejercicio sobre una página de inicio de sesión que hace uso de HTTPS (login.gov), ingresando el usuario de prueba PROBANDO y una contraseña ficticia, con el fin de verificar si esas mismas credenciales podían localizarse en la captura como ocurrió con HTTP.

\begin{lstlisting}
tls
tls.handshake.type == 1
\end{lstlisting}

\captura{https-tls-handshake.png}{Credenciales de prueba ingresadas en el formulario de login.gov, servido sobre HTTPS.}

\begin{analisis}
En la figura se observa el formulario de inicio de sesión con el usuario PROBANDO y una contraseña de prueba, justo antes de enviarlo. A diferencia del caso HTTP, el navegador no envía estos datos directamente: primero establece un canal TLS con el servidor (el handshake, que inicia con el Client Hello), y solo cuando el canal cifrado está listo transmite el formulario, de modo que lo que se escribe en pantalla no es lo que viaja por la red.
\end{analisis}

\captura{https-application-data.png}{Captura del tráfico durante el inicio de sesión sobre HTTPS: las credenciales no aparecen en ningún paquete.}

\begin{analisis}
Al revisar la captura realizada durante el envío del formulario, ordenada por protocolo, solo se observan protocolos de servicio de la red local (ARP, NetBIOS, DHCPv6, DNS, LLMNR y mDNS) y ningún paquete expone el usuario PROBANDO ni la contraseña. El contraste con el caso HTTP es directo: allí bastaba un filtro para aislar el POST y leer el formulario en claro, mientras que aquí las credenciales viajan dentro de registros TLS Application Data, de los que Wireshark solo puede mostrar la cabecera y la longitud del contenido cifrado, sin forma de reconstruirlas sin las claves de sesión.extremos de la comunicación.
\end{analisis}

\pregunta{Al analizar los paquetes HTTPS capturados, ¿es posible ver la información de las credenciales? Justifique su respuesta.}
\begin{respuesta}
No es posible ver las credenciales desde una posición pasiva en la red, porque HTTPS es el mismo HTTP transportado dentro de un canal TLS. Antes de que el navegador envíe un solo byte del formulario se ejecuta el handshake: en el Client Hello y el Server Hello se negocian los parámetros y se acuerdan las claves de sesión, y luego el servidor presenta su certificado. A partir de ese punto todo el intercambio viaja como registros Application Data, de los que Wireshark solo puede mostrar la cabecera y la longitud del contenido cifrado. El usuario y la contraseña solo existen en claro en los extremos (navegador y servidor); la única forma de leerlos sería tener acceso a uno de ellos y exportar las claves de sesión (por ejemplo, con la variable SSLKEYLOGFILE del navegador), que Wireshark puede usar para descifrar la sesión desde Preferences → Protocols → TLS.


El cifrado, sin embargo, no oculta toda la conversación: siguen siendo visibles las direcciones IP y los puertos, la consulta DNS previa, la extensión SNI del Client Hello (salvo que se use Encrypted Client Hello), el certificado del servidor en TLS 1.2 (en TLS 1.3 viaja cifrado) y el volumen y la temporización de los paquetes. Dicho de otra forma, HTTPS protege el contenido de la sesión pero no oculta con quién se habla, lo que en un análisis forense sigue siendo información aprovechable.



\end{respuesta}
